\documentclass[11pt,a4paper]{article}

\def\courseTitle{Industrial Security (Bachelor)}
\def\labNumber{04}
\def\labTitle{Three Protocols on the Wire -- Modbus, OPC UA, MQTT}
\def\labSection{04}
\def\labTime{90--120 min}

% =====================================================================
% Shared preamble for lab sheets.
%
% Caller must \documentclass{article} first, then define:
%   \def\courseTitle{...}
%   \def\labNumber{NN}
%   \def\labTitle{...}
%   \def\labTime{e.g. "30--45 min"}
%   \def\labSection{e.g. "02"}
% Then % =====================================================================
% Shared preamble for lab sheets.
%
% Caller must \documentclass{article} first, then define:
%   \def\courseTitle{...}
%   \def\labNumber{NN}
%   \def\labTitle{...}
%   \def\labTime{e.g. "30--45 min"}
%   \def\labSection{e.g. "02"}
% Then % =====================================================================
% Shared preamble for lab sheets.
%
% Caller must \documentclass{article} first, then define:
%   \def\courseTitle{...}
%   \def\labNumber{NN}
%   \def\labTitle{...}
%   \def\labTime{e.g. "30--45 min"}
%   \def\labSection{e.g. "02"}
% Then \input{../../template/lab-preamble.tex}
% =====================================================================

\usepackage[utf8]{inputenc}
\usepackage[T1]{fontenc}
\usepackage[english]{babel}
\usepackage[a4paper, margin=2cm]{geometry}
\usepackage{graphicx}
\usepackage{listings}
\usepackage{xcolor}
\usepackage{colortbl}
\usepackage{tikz}
\usepackage{amsmath}
\usepackage{amssymb}
\usepackage{booktabs}
\usepackage{enumitem}
\usepackage{titlesec}
\usepackage{fancyhdr}
\usepackage{hyperref}
\usepackage{tcolorbox}
\tcbuselibrary{skins, breakable}
\usepackage{fontawesome5}

\input{../../../template/tikz-styles.tex}

\providecommand{\courseAuthor}{Industrial Security Lecture Series}
\providecommand{\courseInstitute}{Department of Computer Science}
\providecommand{\companyLogo}{logo}

\definecolor{slideAccent}{HTML}{00205B}
\definecolor{slideSecondary}{HTML}{00A0DC}

% --- Corporate branding overrides ------------------------------------
\IfFileExists{../../../branding.tex}{\input{../../../branding.tex}}{}

\colorlet{labAccent}{slideAccent}
\colorlet{labSec}{slideSecondary}
\definecolor{labCheck}{HTML}{2E7D32}
\definecolor{labWarn}{HTML}{B23A48}

\lstset{
  backgroundcolor=\color{black!4},
  basicstyle=\ttfamily\footnotesize,
  breaklines=true,
  frame=leftline,
  framerule=2pt,
  rulecolor=\color{labAccent},
  numbers=left,
  numberstyle=\tiny\color{black!50},
  showstringspaces=false,
  tabsize=2
}

\setlength{\tabcolsep}{4pt}

% Let TeX add a little extra inter-word stretch as a last resort before
% it reports an Overfull \hbox. Only kicks in on lines that would
% otherwise overflow (long \texttt paths, URLs, CIDRs); never loosens a
% line that already fits. Keeps prose/itemize within the margin.
\setlength{\emergencystretch}{3em}

\pagestyle{fancy}
\fancyhf{}
\renewcommand{\headrulewidth}{0.4pt}
\lhead{\small \courseTitle{} -- Lab \labNumber}
\rhead{\small Maps to Section \labSection}
\cfoot{\thepage}

\newtcolorbox{stepbox}[1]{
  enhanced, breakable, arc=2pt, boxrule=0pt,
  colback=labAccent!7, colframe=labAccent, leftrule=3pt,
  fonttitle=\bfseries\small, coltitle=white,
  colbacktitle=labAccent,
  title={\faTools\ Step #1}
}

\newtcolorbox{warnbox}{
  enhanced, breakable, arc=2pt, boxrule=0pt,
  colback=labWarn!7, colframe=labWarn, leftrule=4pt,
  fonttitle=\bfseries\small, coltitle=white, colbacktitle=labWarn,
  title={\faExclamationTriangle\ Caution}
}

\newtcolorbox{notebox}{
  enhanced, breakable, arc=2pt, boxrule=0pt,
  colback=labAccent!4, colframe=labAccent, leftrule=2pt,
  fonttitle=\bfseries\small, coltitle=white, colbacktitle=labAccent,
  title={\faInfoCircle\ Note}
}

\newtcolorbox{goalbox}{
  enhanced, breakable, arc=2pt, boxrule=0pt,
  colback=labCheck!7, colframe=labCheck, leftrule=3pt,
  fonttitle=\bfseries\small, coltitle=white, colbacktitle=labCheck,
  title={\faBullseye\ Goal}
}

\newtcolorbox{deliverbox}{
  enhanced, breakable, arc=2pt, boxrule=0pt,
  colback=labSec!10, colframe=labSec, leftrule=3pt,
  fonttitle=\bfseries\small, coltitle=white, colbacktitle=labSec,
  title={\faClipboardList\ Deliverable}
}

% --- Solution-sheet boxes (used only by lab-solutions.tex) ----------
\newtcolorbox{solutionbox}[1]{
  enhanced, breakable, arc=2pt, boxrule=0pt,
  colback=labCheck!7, colframe=labCheck, leftrule=3pt,
  fonttitle=\bfseries\small, coltitle=white, colbacktitle=labCheck,
  title={\faCheckCircle\ #1}
}

\newtcolorbox{rubricbox}{
  enhanced, breakable, arc=2pt, boxrule=0pt,
  colback=labSec!7, colframe=labSec, leftrule=3pt,
  fonttitle=\bfseries\small, coltitle=white, colbacktitle=labSec,
  title={\faBalanceScale\ Grading rubric}
}

\titleformat{\section}{\large\bfseries\color{labAccent}}{\thesection}{0.5em}{}

\newcommand{\labheader}{%
  \noindent
  \begin{tikzpicture}
    \fill[labAccent] (0,0) rectangle (\textwidth, -1.6cm);
    \node[anchor=north west, text=white, font=\bfseries\large] at (0.6cm, -0.4cm) {Lab \labNumber{} -- \labTitle};
    \node[anchor=south west, text=white, font=\small] at (0.6cm, -1.3cm) {\courseTitle{} \quad $\bullet$ \quad Maps to Section \labSection{} \quad $\bullet$ \quad \labTime};
    \node[anchor=north east, fill=labSec, text=white, font=\bfseries, inner sep=4pt, rounded corners=2pt]
      at ([xshift=-0.4cm, yshift=-0.4cm]\textwidth,0) {LAB \labNumber};
  \end{tikzpicture}
  \vspace{4mm}
}

% =====================================================================

\usepackage[utf8]{inputenc}
\usepackage[T1]{fontenc}
\usepackage[english]{babel}
\usepackage[a4paper, margin=2cm]{geometry}
\usepackage{graphicx}
\usepackage{listings}
\usepackage{xcolor}
\usepackage{colortbl}
\usepackage{tikz}
\usepackage{amsmath}
\usepackage{amssymb}
\usepackage{booktabs}
\usepackage{enumitem}
\usepackage{titlesec}
\usepackage{fancyhdr}
\usepackage{hyperref}
\usepackage{tcolorbox}
\tcbuselibrary{skins, breakable}
\usepackage{fontawesome5}

% =====================================================================
% Shared TikZ styles for the Industrial Security lecture series.
% Loaded by every slide deck and exercise sheet that uses TikZ.
% =====================================================================

\usetikzlibrary{
  arrows.meta,
  shapes,
  shapes.geometric,
  shapes.symbols,
  shapes.multipart,
  positioning,
  fit,
  calc,
  backgrounds,
  decorations.pathmorphing,
  decorations.pathreplacing,
  decorations.text,
  patterns,
  matrix,
  shadows
}

% --- colour palette --------------------------------------------------
\definecolor{isOT}{HTML}{1F4E79}     % OT (deep blue)
\definecolor{isIT}{HTML}{6F2DA8}     % IT (purple)
\definecolor{isSafety}{HTML}{C00000} % safety red
\definecolor{isNet}{HTML}{2E7D32}    % network green
\definecolor{isAttack}{HTML}{B23A48} % attack red
\definecolor{isAccent}{HTML}{E08E0B} % amber
\definecolor{isMute}{HTML}{6B6B6B}   % grey
\definecolor{isLight}{HTML}{EDF2F8}  % light background
\definecolor{isPLC}{HTML}{2C5282}
\definecolor{isHMI}{HTML}{2F855A}
\definecolor{isSCADA}{HTML}{6B46C1}

% --- generic node styles --------------------------------------------
\tikzset{
  isnode/.style={
    draw, rounded corners=2pt, minimum height=8mm, minimum width=18mm,
    font=\small, align=center, thick
  },
  isbox/.style={isnode, fill=isLight},
  plc/.style={isnode, fill=isPLC!15, draw=isPLC, thick},
  hmi/.style={isnode, fill=isHMI!15, draw=isHMI, thick},
  scada/.style={isnode, fill=isSCADA!15, draw=isSCADA, thick},
  sensor/.style={isnode, fill=isNet!10, draw=isNet, minimum height=6mm, minimum width=14mm, font=\scriptsize},
  actuator/.style={isnode, fill=isAccent!15, draw=isAccent, minimum height=6mm, minimum width=14mm, font=\scriptsize},
  firewall/.style={isnode, fill=isAttack!15, draw=isAttack, thick},
  server/.style={isnode, fill=isMute!15, draw=isMute!80, thick},
  switch/.style={isnode, fill=isLight, draw=black!60, minimum height=6mm, minimum width=14mm, font=\scriptsize},
  workstation/.style={isnode, fill=white, draw=isOT, thick},
  attacker/.style={isnode, fill=isAttack!20, draw=isAttack, thick, font=\small\bfseries},
  inetnode/.style={draw, ellipse, minimum width=24mm, minimum height=10mm, font=\scriptsize, fill=isLight, thick},
  process/.style={isnode, fill=isOT!15, draw=isOT, thick, minimum width=24mm},
  zone/.style={
    draw, rounded corners=4pt, thick, dashed, inner sep=4mm
  },
  conduit/.style={-{Stealth[length=2.5mm]}, thick, isNet!80!black},
  attack/.style={-{Stealth[length=2.5mm]}, thick, isAttack, dashed},
  flow/.style={-{Stealth[length=2.5mm]}, thick, isOT!70!black},
  netline/.style={thick, black!60},
  axisline/.style={-{Stealth[length=2mm]}, thick, black!70},
  tag/.style={font=\scriptsize\itshape, fill=white, inner sep=1pt},
  level/.style={
    draw, fill=isLight, minimum width=0.85\linewidth, minimum height=8mm,
    font=\small, anchor=west
  },
  brace/.style={decorate, decoration={brace, amplitude=4pt, raise=2pt}},
  legendbox/.style={draw, rounded corners=2pt, fill=isLight, inner sep=2mm, font=\scriptsize, align=left}
}

% --- handy macros ----------------------------------------------------
% \plcicon, \hmiicon etc as simple labelled boxes; useful inline.
\newcommand{\plcicon}[1][PLC]{\tikz[baseline=-0.6ex]{\node[plc, minimum width=10mm, minimum height=5mm, font=\scriptsize, inner sep=1pt]{#1};}}
\newcommand{\hmiicon}[1][HMI]{\tikz[baseline=-0.6ex]{\node[hmi, minimum width=10mm, minimum height=5mm, font=\scriptsize, inner sep=1pt]{#1};}}
\newcommand{\fwicon}[1][FW]{\tikz[baseline=-0.6ex]{\node[firewall, minimum width=10mm, minimum height=5mm, font=\scriptsize, inner sep=1pt]{#1};}}


\providecommand{\courseAuthor}{Industrial Security Lecture Series}
\providecommand{\courseInstitute}{Department of Computer Science}
\providecommand{\companyLogo}{logo}

\definecolor{slideAccent}{HTML}{00205B}
\definecolor{slideSecondary}{HTML}{00A0DC}

% --- Corporate branding overrides ------------------------------------
\IfFileExists{../../../branding.tex}{% =====================================================================
% Corporate / institutional branding -- single file to customise the
% look of the entire course (slides, notes, exercises, labs).
%
% HOW TO REBRAND IN UNDER A MINUTE:
%   1. Replace `branding/logo.pdf` with your own logo
%      (PDF preferred; PNG and JPG also work -- name it `logo.<ext>`).
%   2. (Optional) change the two colours below to your brand palette.
%   3. (Optional) override the author/institute lines below.
%   4. Re-run `make` at the repo root. That's it.
%
% Everything in this file has sensible defaults; you can delete a line
% to fall back to the built-in defaults, or comment it out.
%
% This file is loaded automatically by the slides, exercise, and lab
% preambles -- you never need to \input it manually.
% =====================================================================

% --- Logo --------------------------------------------------------------
% Path is resolved from each leaf-build directory; the default points at
% the shared `branding/` folder at the repo root. If you put your logo
% somewhere else, set the full or relative path here (without extension):
\renewcommand{\companyLogo}{../../../branding/logo}

% --- Author / institute (appears on the title page footer) ------------
\renewcommand{\courseAuthor}{}
\renewcommand{\courseInstitute}{}

% --- Brand colours ----------------------------------------------------
% slideAccent      = primary brand colour (title bands, accent rule)
% slideSecondary   = secondary brand colour (section badge, highlight)
% Both use HTML hex codes. Keep enough contrast against white text.
\definecolor{slideAccent}{HTML}{00205B}      % deep navy
\definecolor{slideSecondary}{HTML}{00A0DC}   % light blue accent
}{}

\colorlet{labAccent}{slideAccent}
\colorlet{labSec}{slideSecondary}
\definecolor{labCheck}{HTML}{2E7D32}
\definecolor{labWarn}{HTML}{B23A48}

\lstset{
  backgroundcolor=\color{black!4},
  basicstyle=\ttfamily\footnotesize,
  breaklines=true,
  frame=leftline,
  framerule=2pt,
  rulecolor=\color{labAccent},
  numbers=left,
  numberstyle=\tiny\color{black!50},
  showstringspaces=false,
  tabsize=2
}

\setlength{\tabcolsep}{4pt}

% Let TeX add a little extra inter-word stretch as a last resort before
% it reports an Overfull \hbox. Only kicks in on lines that would
% otherwise overflow (long \texttt paths, URLs, CIDRs); never loosens a
% line that already fits. Keeps prose/itemize within the margin.
\setlength{\emergencystretch}{3em}

\pagestyle{fancy}
\fancyhf{}
\renewcommand{\headrulewidth}{0.4pt}
\lhead{\small \courseTitle{} -- Lab \labNumber}
\rhead{\small Maps to Section \labSection}
\cfoot{\thepage}

\newtcolorbox{stepbox}[1]{
  enhanced, breakable, arc=2pt, boxrule=0pt,
  colback=labAccent!7, colframe=labAccent, leftrule=3pt,
  fonttitle=\bfseries\small, coltitle=white,
  colbacktitle=labAccent,
  title={\faTools\ Step #1}
}

\newtcolorbox{warnbox}{
  enhanced, breakable, arc=2pt, boxrule=0pt,
  colback=labWarn!7, colframe=labWarn, leftrule=4pt,
  fonttitle=\bfseries\small, coltitle=white, colbacktitle=labWarn,
  title={\faExclamationTriangle\ Caution}
}

\newtcolorbox{notebox}{
  enhanced, breakable, arc=2pt, boxrule=0pt,
  colback=labAccent!4, colframe=labAccent, leftrule=2pt,
  fonttitle=\bfseries\small, coltitle=white, colbacktitle=labAccent,
  title={\faInfoCircle\ Note}
}

\newtcolorbox{goalbox}{
  enhanced, breakable, arc=2pt, boxrule=0pt,
  colback=labCheck!7, colframe=labCheck, leftrule=3pt,
  fonttitle=\bfseries\small, coltitle=white, colbacktitle=labCheck,
  title={\faBullseye\ Goal}
}

\newtcolorbox{deliverbox}{
  enhanced, breakable, arc=2pt, boxrule=0pt,
  colback=labSec!10, colframe=labSec, leftrule=3pt,
  fonttitle=\bfseries\small, coltitle=white, colbacktitle=labSec,
  title={\faClipboardList\ Deliverable}
}

% --- Solution-sheet boxes (used only by lab-solutions.tex) ----------
\newtcolorbox{solutionbox}[1]{
  enhanced, breakable, arc=2pt, boxrule=0pt,
  colback=labCheck!7, colframe=labCheck, leftrule=3pt,
  fonttitle=\bfseries\small, coltitle=white, colbacktitle=labCheck,
  title={\faCheckCircle\ #1}
}

\newtcolorbox{rubricbox}{
  enhanced, breakable, arc=2pt, boxrule=0pt,
  colback=labSec!7, colframe=labSec, leftrule=3pt,
  fonttitle=\bfseries\small, coltitle=white, colbacktitle=labSec,
  title={\faBalanceScale\ Grading rubric}
}

\titleformat{\section}{\large\bfseries\color{labAccent}}{\thesection}{0.5em}{}

\newcommand{\labheader}{%
  \noindent
  \begin{tikzpicture}
    \fill[labAccent] (0,0) rectangle (\textwidth, -1.6cm);
    \node[anchor=north west, text=white, font=\bfseries\large] at (0.6cm, -0.4cm) {Lab \labNumber{} -- \labTitle};
    \node[anchor=south west, text=white, font=\small] at (0.6cm, -1.3cm) {\courseTitle{} \quad $\bullet$ \quad Maps to Section \labSection{} \quad $\bullet$ \quad \labTime};
    \node[anchor=north east, fill=labSec, text=white, font=\bfseries, inner sep=4pt, rounded corners=2pt]
      at ([xshift=-0.4cm, yshift=-0.4cm]\textwidth,0) {LAB \labNumber};
  \end{tikzpicture}
  \vspace{4mm}
}

% =====================================================================

\usepackage[utf8]{inputenc}
\usepackage[T1]{fontenc}
\usepackage[english]{babel}
\usepackage[a4paper, margin=2cm]{geometry}
\usepackage{graphicx}
\usepackage{listings}
\usepackage{xcolor}
\usepackage{colortbl}
\usepackage{tikz}
\usepackage{amsmath}
\usepackage{amssymb}
\usepackage{booktabs}
\usepackage{enumitem}
\usepackage{titlesec}
\usepackage{fancyhdr}
\usepackage{hyperref}
\usepackage{tcolorbox}
\tcbuselibrary{skins, breakable}
\usepackage{fontawesome5}

% =====================================================================
% Shared TikZ styles for the Industrial Security lecture series.
% Loaded by every slide deck and exercise sheet that uses TikZ.
% =====================================================================

\usetikzlibrary{
  arrows.meta,
  shapes,
  shapes.geometric,
  shapes.symbols,
  shapes.multipart,
  positioning,
  fit,
  calc,
  backgrounds,
  decorations.pathmorphing,
  decorations.pathreplacing,
  decorations.text,
  patterns,
  matrix,
  shadows
}

% --- colour palette --------------------------------------------------
\definecolor{isOT}{HTML}{1F4E79}     % OT (deep blue)
\definecolor{isIT}{HTML}{6F2DA8}     % IT (purple)
\definecolor{isSafety}{HTML}{C00000} % safety red
\definecolor{isNet}{HTML}{2E7D32}    % network green
\definecolor{isAttack}{HTML}{B23A48} % attack red
\definecolor{isAccent}{HTML}{E08E0B} % amber
\definecolor{isMute}{HTML}{6B6B6B}   % grey
\definecolor{isLight}{HTML}{EDF2F8}  % light background
\definecolor{isPLC}{HTML}{2C5282}
\definecolor{isHMI}{HTML}{2F855A}
\definecolor{isSCADA}{HTML}{6B46C1}

% --- generic node styles --------------------------------------------
\tikzset{
  isnode/.style={
    draw, rounded corners=2pt, minimum height=8mm, minimum width=18mm,
    font=\small, align=center, thick
  },
  isbox/.style={isnode, fill=isLight},
  plc/.style={isnode, fill=isPLC!15, draw=isPLC, thick},
  hmi/.style={isnode, fill=isHMI!15, draw=isHMI, thick},
  scada/.style={isnode, fill=isSCADA!15, draw=isSCADA, thick},
  sensor/.style={isnode, fill=isNet!10, draw=isNet, minimum height=6mm, minimum width=14mm, font=\scriptsize},
  actuator/.style={isnode, fill=isAccent!15, draw=isAccent, minimum height=6mm, minimum width=14mm, font=\scriptsize},
  firewall/.style={isnode, fill=isAttack!15, draw=isAttack, thick},
  server/.style={isnode, fill=isMute!15, draw=isMute!80, thick},
  switch/.style={isnode, fill=isLight, draw=black!60, minimum height=6mm, minimum width=14mm, font=\scriptsize},
  workstation/.style={isnode, fill=white, draw=isOT, thick},
  attacker/.style={isnode, fill=isAttack!20, draw=isAttack, thick, font=\small\bfseries},
  inetnode/.style={draw, ellipse, minimum width=24mm, minimum height=10mm, font=\scriptsize, fill=isLight, thick},
  process/.style={isnode, fill=isOT!15, draw=isOT, thick, minimum width=24mm},
  zone/.style={
    draw, rounded corners=4pt, thick, dashed, inner sep=4mm
  },
  conduit/.style={-{Stealth[length=2.5mm]}, thick, isNet!80!black},
  attack/.style={-{Stealth[length=2.5mm]}, thick, isAttack, dashed},
  flow/.style={-{Stealth[length=2.5mm]}, thick, isOT!70!black},
  netline/.style={thick, black!60},
  axisline/.style={-{Stealth[length=2mm]}, thick, black!70},
  tag/.style={font=\scriptsize\itshape, fill=white, inner sep=1pt},
  level/.style={
    draw, fill=isLight, minimum width=0.85\linewidth, minimum height=8mm,
    font=\small, anchor=west
  },
  brace/.style={decorate, decoration={brace, amplitude=4pt, raise=2pt}},
  legendbox/.style={draw, rounded corners=2pt, fill=isLight, inner sep=2mm, font=\scriptsize, align=left}
}

% --- handy macros ----------------------------------------------------
% \plcicon, \hmiicon etc as simple labelled boxes; useful inline.
\newcommand{\plcicon}[1][PLC]{\tikz[baseline=-0.6ex]{\node[plc, minimum width=10mm, minimum height=5mm, font=\scriptsize, inner sep=1pt]{#1};}}
\newcommand{\hmiicon}[1][HMI]{\tikz[baseline=-0.6ex]{\node[hmi, minimum width=10mm, minimum height=5mm, font=\scriptsize, inner sep=1pt]{#1};}}
\newcommand{\fwicon}[1][FW]{\tikz[baseline=-0.6ex]{\node[firewall, minimum width=10mm, minimum height=5mm, font=\scriptsize, inner sep=1pt]{#1};}}


\providecommand{\courseAuthor}{Industrial Security Lecture Series}
\providecommand{\courseInstitute}{Department of Computer Science}
\providecommand{\companyLogo}{logo}

\definecolor{slideAccent}{HTML}{00205B}
\definecolor{slideSecondary}{HTML}{00A0DC}

% --- Corporate branding overrides ------------------------------------
\IfFileExists{../../../branding.tex}{% =====================================================================
% Corporate / institutional branding -- single file to customise the
% look of the entire course (slides, notes, exercises, labs).
%
% HOW TO REBRAND IN UNDER A MINUTE:
%   1. Replace `branding/logo.pdf` with your own logo
%      (PDF preferred; PNG and JPG also work -- name it `logo.<ext>`).
%   2. (Optional) change the two colours below to your brand palette.
%   3. (Optional) override the author/institute lines below.
%   4. Re-run `make` at the repo root. That's it.
%
% Everything in this file has sensible defaults; you can delete a line
% to fall back to the built-in defaults, or comment it out.
%
% This file is loaded automatically by the slides, exercise, and lab
% preambles -- you never need to \input it manually.
% =====================================================================

% --- Logo --------------------------------------------------------------
% Path is resolved from each leaf-build directory; the default points at
% the shared `branding/` folder at the repo root. If you put your logo
% somewhere else, set the full or relative path here (without extension):
\renewcommand{\companyLogo}{../../../branding/logo}

% --- Author / institute (appears on the title page footer) ------------
\renewcommand{\courseAuthor}{}
\renewcommand{\courseInstitute}{}

% --- Brand colours ----------------------------------------------------
% slideAccent      = primary brand colour (title bands, accent rule)
% slideSecondary   = secondary brand colour (section badge, highlight)
% Both use HTML hex codes. Keep enough contrast against white text.
\definecolor{slideAccent}{HTML}{00205B}      % deep navy
\definecolor{slideSecondary}{HTML}{00A0DC}   % light blue accent
}{}

\colorlet{labAccent}{slideAccent}
\colorlet{labSec}{slideSecondary}
\definecolor{labCheck}{HTML}{2E7D32}
\definecolor{labWarn}{HTML}{B23A48}

\lstset{
  backgroundcolor=\color{black!4},
  basicstyle=\ttfamily\footnotesize,
  breaklines=true,
  frame=leftline,
  framerule=2pt,
  rulecolor=\color{labAccent},
  numbers=left,
  numberstyle=\tiny\color{black!50},
  showstringspaces=false,
  tabsize=2
}

\setlength{\tabcolsep}{4pt}

% Let TeX add a little extra inter-word stretch as a last resort before
% it reports an Overfull \hbox. Only kicks in on lines that would
% otherwise overflow (long \texttt paths, URLs, CIDRs); never loosens a
% line that already fits. Keeps prose/itemize within the margin.
\setlength{\emergencystretch}{3em}

\pagestyle{fancy}
\fancyhf{}
\renewcommand{\headrulewidth}{0.4pt}
\lhead{\small \courseTitle{} -- Lab \labNumber}
\rhead{\small Maps to Section \labSection}
\cfoot{\thepage}

\newtcolorbox{stepbox}[1]{
  enhanced, breakable, arc=2pt, boxrule=0pt,
  colback=labAccent!7, colframe=labAccent, leftrule=3pt,
  fonttitle=\bfseries\small, coltitle=white,
  colbacktitle=labAccent,
  title={\faTools\ Step #1}
}

\newtcolorbox{warnbox}{
  enhanced, breakable, arc=2pt, boxrule=0pt,
  colback=labWarn!7, colframe=labWarn, leftrule=4pt,
  fonttitle=\bfseries\small, coltitle=white, colbacktitle=labWarn,
  title={\faExclamationTriangle\ Caution}
}

\newtcolorbox{notebox}{
  enhanced, breakable, arc=2pt, boxrule=0pt,
  colback=labAccent!4, colframe=labAccent, leftrule=2pt,
  fonttitle=\bfseries\small, coltitle=white, colbacktitle=labAccent,
  title={\faInfoCircle\ Note}
}

\newtcolorbox{goalbox}{
  enhanced, breakable, arc=2pt, boxrule=0pt,
  colback=labCheck!7, colframe=labCheck, leftrule=3pt,
  fonttitle=\bfseries\small, coltitle=white, colbacktitle=labCheck,
  title={\faBullseye\ Goal}
}

\newtcolorbox{deliverbox}{
  enhanced, breakable, arc=2pt, boxrule=0pt,
  colback=labSec!10, colframe=labSec, leftrule=3pt,
  fonttitle=\bfseries\small, coltitle=white, colbacktitle=labSec,
  title={\faClipboardList\ Deliverable}
}

% --- Solution-sheet boxes (used only by lab-solutions.tex) ----------
\newtcolorbox{solutionbox}[1]{
  enhanced, breakable, arc=2pt, boxrule=0pt,
  colback=labCheck!7, colframe=labCheck, leftrule=3pt,
  fonttitle=\bfseries\small, coltitle=white, colbacktitle=labCheck,
  title={\faCheckCircle\ #1}
}

\newtcolorbox{rubricbox}{
  enhanced, breakable, arc=2pt, boxrule=0pt,
  colback=labSec!7, colframe=labSec, leftrule=3pt,
  fonttitle=\bfseries\small, coltitle=white, colbacktitle=labSec,
  title={\faBalanceScale\ Grading rubric}
}

\titleformat{\section}{\large\bfseries\color{labAccent}}{\thesection}{0.5em}{}

\newcommand{\labheader}{%
  \noindent
  \begin{tikzpicture}
    \fill[labAccent] (0,0) rectangle (\textwidth, -1.6cm);
    \node[anchor=north west, text=white, font=\bfseries\large] at (0.6cm, -0.4cm) {Lab \labNumber{} -- \labTitle};
    \node[anchor=south west, text=white, font=\small] at (0.6cm, -1.3cm) {\courseTitle{} \quad $\bullet$ \quad Maps to Section \labSection{} \quad $\bullet$ \quad \labTime};
    \node[anchor=north east, fill=labSec, text=white, font=\bfseries, inner sep=4pt, rounded corners=2pt]
      at ([xshift=-0.4cm, yshift=-0.4cm]\textwidth,0) {LAB \labNumber};
  \end{tikzpicture}
  \vspace{4mm}
}


\begin{document}
\labheader

\begin{goalbox}
Touch three industrial protocols in the same lab session and watch the
security delta with your own eyes. Read and write Modbus without
authentication, browse an OPC UA server in anonymous mode, publish and
wildcard-subscribe an MQTT topic, then fill in a side-by-side comparison
table. Stretch goal: spin up a fake PLC in 20 lines of Python.
\end{goalbox}

\section*{0 -- Mental model before you start}

Three IT analogues will carry you through the whole lab. Keep them in mind
as you run each step:

\begin{itemize}[leftmargin=*, nosep]
  \item \textbf{Modbus} $\approx$ a tiny key--value store. Sixteen-bit
        addresses, no schema, no auth. Think Redis on a public IP with
        no \texttt{requirepass}.
  \item \textbf{OPC UA} $\approx$ a typed REST resource tree. Browse
        nodes, read attributes, subscribe to changes. Think a JSON API
        with optional mTLS -- where the lab server has the TLS turned
        off.
  \item \textbf{MQTT} $\approx$ a pub/sub broker. Think Kafka or Redis
        pub/sub. The broker is the choke point: own it and you own
        everything that flows through it.
\end{itemize}

You do not need to know any electrical engineering or ladder logic to
finish this lab. Every ``register'', ``coil'' and ``node'' you touch is
just a memory address with a different label on it.

\section*{1 -- Pre-flight}

\begin{stepbox}{1.1 -- Bring the stack up}
From the repository root:
\begin{lstlisting}[language=bash]
cd bachelor/labs/_stack
docker compose up -d
docker compose ps
\end{lstlisting}
\emph{Expected:} services \texttt{openplc}, \texttt{opcua}, \texttt{mqtt},
and \texttt{student} report \texttt{running} (and \texttt{openplc-init}
sits in \texttt{exited (0)} once it has uploaded the demo program).
\end{stepbox}

\begin{stepbox}{1.2 -- Verify each service answers}
\begin{lstlisting}[language=bash]
docker compose exec student bash -c '
  echo "--- Modbus (502/tcp)";   nc -zv 172.28.0.10 502  || true
  echo "--- OPC UA (4840/tcp)";  nc -zv 172.28.0.11 4840 || true
  echo "--- MQTT   (1883/tcp)";  nc -zv 172.28.0.12 1883 || true
'
\end{lstlisting}
\emph{Expected:} three lines saying \texttt{open} or \texttt{succeeded}.
If any line says \texttt{refused}, re-check \texttt{docker compose ps}
before continuing -- nothing in the lab works without all three.
\end{stepbox}

\begin{stepbox}{1.3 -- Open a working shell in the student container}
\begin{lstlisting}[language=bash]
docker compose exec student bash
\end{lstlisting}
Everything from here on runs \emph{inside} that container unless the step
says otherwise. The container has \texttt{python3}, \texttt{pymodbus},
\texttt{asyncua}, \texttt{paho-mqtt}, \texttt{mosquitto\_pub}/
\texttt{mosquitto\_sub}, \texttt{tcpdump}, and \texttt{tshark} pre-installed.
\end{stepbox}

\section*{2 -- Modbus: a key--value store with no password}

\begin{stepbox}{2.1 -- Honest read of holding registers}
Holding registers are 16-bit cells indexed 0..65535. Think of them as
\texttt{GET key} on a Redis instance with integer keys.
\begin{lstlisting}[language=bash]
python3 /labs/scripts/modbus_read.py
\end{lstlisting}
\emph{Expected output} (freshly booted PLC):
\begin{lstlisting}
Registers 0..9: [0, 0, 0, 0, 0, 0, 0, 0, 0, 0]
\end{lstlisting}
Non-zero values just mean a previous lab run left state behind --
\texttt{docker compose restart openplc} resets them.
\end{stepbox}

\begin{stepbox}{2.2 -- Capture the request on the wire}
Start a packet capture, do one read, then stop the capture:
\begin{lstlisting}[language=bash]
tcpdump -i eth0 -w /labs/captures/lab04-mb.pcap port 502 &
TCPDUMP_PID=$!
sleep 1
python3 /labs/scripts/modbus_read.py
kill -INT $TCPDUMP_PID
wait $TCPDUMP_PID 2>/dev/null
ls -lh /labs/captures/lab04-mb.pcap
\end{lstlisting}
\emph{Expected:} a non-empty pcap file, typically 1--3\,KB.
\end{stepbox}

\begin{stepbox}{2.3 -- Decode the MBAP header byte by byte}
Modbus TCP wraps the original serial PDU in a 7-byte \emph{MBAP} header.
Use \texttt{tshark} to print one full request:
\begin{lstlisting}[language=bash]
tshark -r /labs/captures/lab04-mb.pcap -V -Y modbus | head -60
\end{lstlisting}
\emph{Expected fields} in the request:
\begin{lstlisting}
Modbus/TCP
    Transaction Identifier: 1
    Protocol Identifier:    0
    Length:                 6
    Unit Identifier:        1
Modbus
    .000 0011 = Function Code: Read Holding Registers (3)
    Reference Number:       0
    Word Count:             10
\end{lstlisting}
Now line that up with the seven MBAP bytes plus the PDU:
\begin{center}\small
\begin{tabular}{@{}lll@{}}
\toprule
\textbf{Bytes} & \textbf{Field} & \textbf{IT analogue} \\
\midrule
\texttt{00 01} & Transaction ID  & request correlation cookie \\
\texttt{00 00} & Protocol ID     & always 0 for Modbus \\
\texttt{00 06} & Length          & bytes that follow \\
\texttt{01}    & Unit ID         & gateway / slave selector \\
\texttt{03}    & Function code   & verb -- here ``read regs'' \\
\texttt{00 00} & Reference no.   & start key \\
\texttt{00 0A} & Word count      & number of keys to read \\
\bottomrule
\end{tabular}
\end{center}
There is no token, no MAC, no session ID. Anyone who can send these
twelve bytes is ``authenticated''.
\end{stepbox}

\begin{stepbox}{2.4 -- Write without authentication}
Function code \texttt{0x05} is \emph{Write Single Coil}. A coil is a
1-bit output on the PLC -- in the demo program, coil \texttt{0} drives
output \texttt{QX0.0}.
\begin{lstlisting}[language=bash]
python3 /labs/scripts/modbus_inject.py
\end{lstlisting}
\emph{Expected output:}
\begin{lstlisting}
Write result: WriteSingleCoilResponse(0, on)
\end{lstlisting}
Open the OpenPLC web UI at \texttt{http://127.0.0.1:8080} (default login
\texttt{openplc} / \texttt{openplc}) and check the \emph{Monitoring}
page: \texttt{QX0.0} is now \texttt{TRUE}. Nothing logged the change,
nothing asked who you were.
\end{stepbox}

\begin{stepbox}{2.5 -- Confirm by reading back}
\begin{lstlisting}[language=bash]
python3 -c '
from pymodbus.client import ModbusTcpClient
c = ModbusTcpClient("172.28.0.10", port=502); c.connect()
print("Coil 0:", c.read_coils(address=0, count=1, slave=1).bits[0])
c.close()
'
\end{lstlisting}
\emph{Expected:} \texttt{Coil 0: True}. The bit you just flipped is now
the ground truth of the ``physical'' output -- with zero credentials.
\end{stepbox}

\begin{warnbox}
This works because Modbus has \emph{no} authentication, not because of a
bug. Never aim these scripts at a device you do not own. The same six
bytes that flipped a simulated coil can trip a real breaker in a real
substation.
\end{warnbox}

\section*{3 -- OPC UA: typed REST tree, but anonymous}

\begin{stepbox}{3.1 -- Browse the namespace}
OPC UA exposes a graph of \emph{nodes} the same way a REST API exposes a
tree of resources. The script walks the \texttt{Objects} folder, one
level deep:
\begin{lstlisting}[language=bash]
python3 /labs/scripts/opcua_browse.py
\end{lstlisting}
\emph{Expected output} (open62541 example server):
\begin{lstlisting}
Connected to: opc.tcp://172.28.0.11:4840
 - 0:Server
 - 1:the.answer
 - 1:current-time
 - 1:Demo
\end{lstlisting}
Node names vary slightly across open62541 versions, but \texttt{0:Server}
(the standard root) and at least one user-namespace node
(\texttt{ns=1:...}) must show up.
\end{stepbox}

\begin{stepbox}{3.2 -- Capture the handshake}
\begin{lstlisting}[language=bash]
tcpdump -i eth0 -w /labs/captures/lab04-ua.pcap port 4840 &
TCPDUMP_PID=$!
sleep 1
python3 /labs/scripts/opcua_browse.py
kill -INT $TCPDUMP_PID
wait $TCPDUMP_PID 2>/dev/null
ls -lh /labs/captures/lab04-ua.pcap
\end{lstlisting}
\emph{Expected:} a 4--15\,KB pcap.
\end{stepbox}

\begin{stepbox}{3.3 -- Read the security policy off the wire}
OPC UA has its own dissector in Wireshark / tshark, but you have to tell
it the port:
\begin{lstlisting}[language=bash]
tshark -r /labs/captures/lab04-ua.pcap \
       -d tcp.port==4840,opcua -Y opcua -V | less
\end{lstlisting}
Look for the \texttt{OPN} (OpenSecureChannel) frames. The critical field
is \texttt{SecurityPolicyURI}:
\begin{lstlisting}
SecurityPolicyUri:
  http://opcfoundation.org/UA/SecurityPolicy#None
UserIdentityToken: Anonymous
\end{lstlisting}
\emph{Translation:} no signing, no encryption, no client cert, no user
name. The handshake exists, but it negotiates the empty cipher suite.
Every \texttt{MSG} frame after the OPN -- including the browse results
you just printed -- is plaintext.
\end{stepbox}

\begin{stepbox}{3.4 -- Why the default is ``None''}
Why does the upstream open62541 example ship with
\texttt{SecurityPolicy=None}? Three honest reasons that students should
internalise:
\begin{itemize}[leftmargin=*, nosep]
  \item \textbf{Bootstrapping.} Mutual X.509 needs a CA, a server cert,
        and a client cert. None of that exists out of the box.
  \item \textbf{CPU budget.} A 200\,MHz PLC CPU can spend 10\,\% of its
        scan cycle on one RSA-2048 handshake. That is enough to miss a
        watchdog and trip a process.
  \item \textbf{Interoperability.} Integrators routinely downgrade to
        \texttt{None} to make a third-party HMI talk to a new server.
\end{itemize}
The protocol \emph{can} be secured (\texttt{Basic256Sha256},
\texttt{Aes256\_Sha256\_RsaPss}, signed-and-encrypted, certificate-based
user tokens). The lab server shows you what most field deployments
actually look like.
\end{stepbox}

\section*{4 -- MQTT: Kafka without auth}

The lab broker (\texttt{eclipse-mosquitto}) is configured with
\texttt{allow\_anonymous true} and no ACLs. That is a faithful copy of
the majority of internet-facing MQTT brokers on Shodan today.

\begin{stepbox}{4.1 -- Subscribe to a topic, then publish into it}
Open a \emph{second} student shell (run \texttt{docker compose exec
student bash} again in another terminal). In one shell, subscribe:
\begin{lstlisting}[language=bash]
mosquitto_sub -h 172.28.0.12 -t 'sensors/vibration' -v
\end{lstlisting}
In the other shell, publish:
\begin{lstlisting}[language=bash]
mosquitto_pub -h 172.28.0.12 -t 'sensors/vibration' \
              -m '{"v": 3.7, "ts": 1700000000}'
\end{lstlisting}
\emph{Expected} in the subscriber shell:
\begin{lstlisting}
sensors/vibration {"v": 3.7, "ts": 1700000000}
\end{lstlisting}
You never authenticated. The broker did not care.
\end{stepbox}

\begin{stepbox}{4.2 -- Publish a burst with the helper script}
\begin{lstlisting}[language=bash]
python3 /labs/scripts/mqtt_publish.py
\end{lstlisting}
\emph{Expected:} five JSON messages arrive in the subscriber shell, one
every 200\,ms, the script prints \texttt{published 5 messages...} and
exits.
\end{stepbox}

\begin{stepbox}{4.3 -- Capture and decode an MQTT exchange}
\begin{lstlisting}[language=bash]
tcpdump -i eth0 -w /labs/captures/lab04-mq.pcap port 1883 &
TCPDUMP_PID=$!
sleep 1
python3 /labs/scripts/mqtt_publish.py
kill -INT $TCPDUMP_PID
wait $TCPDUMP_PID 2>/dev/null
tshark -r /labs/captures/lab04-mq.pcap -Y mqtt -V | head -40
\end{lstlisting}
\emph{Expected:} you should see \texttt{CONNECT} (no \texttt{Username}
field), one \texttt{CONNACK} with return code \texttt{0 -- accepted},
then \texttt{PUBLISH} frames with the topic and JSON payload entirely in
the clear. Body and headers are visible to anyone on the segment.
\end{stepbox}

\begin{stepbox}{4.4 -- Wildcard subscription = full telemetry firehose}
MQTT topics are slash-separated paths; \texttt{\#} matches everything
below a point in the tree, \texttt{+} matches a single level.

\emph{Subscribe to the entire tree} (one shell):
\begin{lstlisting}[language=bash]
mosquitto_sub -h 172.28.0.12 -t '#' -v
\end{lstlisting}
\emph{Publish a few topics from the second shell}:
\begin{lstlisting}[language=bash]
mosquitto_pub -h 172.28.0.12 -t 'plant/lineA/pump1/rpm'  -m '1480'
mosquitto_pub -h 172.28.0.12 -t 'plant/lineA/pump1/temp' -m '64'
mosquitto_pub -h 172.28.0.12 -t 'plant/lineB/valve3/pos' -m 'OPEN'
\end{lstlisting}
\emph{Expected in the wildcard subscriber:}
\begin{lstlisting}
plant/lineA/pump1/rpm  1480
plant/lineA/pump1/temp 64
plant/lineB/valve3/pos OPEN
\end{lstlisting}
On a real broker without ACLs that single \texttt{\#} subscription is
the entire plant's telemetry feed, in real time, in plaintext.
\end{stepbox}

\begin{stepbox}{4.5 -- Inject a forged sensor reading}
Nothing stops an unauthenticated client from \emph{publishing} either:
\begin{lstlisting}[language=bash]
mosquitto_pub -h 172.28.0.12 -t 'plant/lineA/pump1/rpm' -m '99999'
\end{lstlisting}
The wildcard subscriber sees \texttt{plant/lineA/pump1/rpm 99999}
indistinguishable from a real measurement. If a downstream dashboard
trusts the broker, it just rendered a false alarm. If a control loop
trusts the broker, it just acted on a lie.
\end{stepbox}

\section*{5 -- The conceptual climax: side-by-side comparison}

\begin{stepbox}{5.1 -- Fill in the table from what you observed}
Compare the three protocols on five dimensions. Treat the lab observations
as evidence, not opinion.

\renewcommand{\arraystretch}{1.25}
\begin{center}\footnotesize
\begin{tabular}{@{}p{2.4cm}p{3.3cm}p{3.3cm}p{3.3cm}@{}}
\toprule
\textbf{Dimension} & \textbf{Modbus TCP} & \textbf{OPC UA (default)} & \textbf{MQTT (plaintext)} \\
\midrule
Authentication
  & None -- 7-byte MBAP, no token
  & Anonymous endpoint, \texttt{SecurityPolicy=None}
  & \texttt{allow\_anonymous true}, no ACL \\
Integrity
  & No MAC, no checksum beyond TCP
  & Sign mode disabled (None policy)
  & TCP only -- no app-layer MAC \\
Confidentiality
  & Cleartext on 502/tcp
  & Cleartext after OPN
  & Cleartext on 1883/tcp \\
Ease of misuse
  & 12 bytes to read or write
  & Browse + write in 10 lines of Python
  & One \texttt{mosquitto\_pub} command \\
Where you'd accept it
  & Air-gapped cell, behind a data-diode or firewall
  & Configured with \texttt{Basic256Sha256}, mutual X.509, no anonymous endpoint
  & TLS on 8883 + per-client ACLs + per-topic permissions \\
\bottomrule
\end{tabular}
\end{center}
\end{stepbox}

\begin{stepbox}{5.2 -- IT analogues, side by side}
Final sanity check -- if you can complete this table from memory, you
have the model:

\renewcommand{\arraystretch}{1.2}
\begin{center}\footnotesize
\begin{tabular}{@{}lll@{}}
\toprule
\textbf{OT concept} & \textbf{IT analogue} & \textbf{What goes wrong} \\
\midrule
Holding register   & Redis integer key      & \texttt{requirepass} not set \\
Coil               & Boolean flag in KV     & Anyone can flip it \\
OPC UA node        & REST resource          & API with no TLS, no auth \\
SecurityPolicy=None & TLS-disabled HTTPS    & Hostname says https, wire does not \\
MQTT topic         & Kafka topic            & No ACL on producer or consumer \\
MQTT \texttt{\#}   & \texttt{SELECT *}      & One subscription = full firehose \\
\bottomrule
\end{tabular}
\end{center}
\end{stepbox}

\section*{6 -- Stretch: spoof a PLC in 20 lines}

\begin{stepbox}{6.1 -- Why this matters}
Reading and writing Modbus is the obvious attack. The subtler one is
\emph{impersonation}: spin up your own server, serve fake values, watch
an HMI or historian accept them as ground truth. Because Modbus has no
identity, the SCADA cannot tell your server apart from the real PLC --
all it sees is ``something on port 502 answered with the right
function code''.
\end{stepbox}

\begin{stepbox}{6.2 -- Write the fake server}
Stop the real PLC's port so the fake one can bind:
\begin{lstlisting}[language=bash]
docker compose stop openplc
\end{lstlisting}
Save as \texttt{/tmp/fake\_plc.py} inside the student container
(\texttt{vim} or \texttt{cat > /tmp/fake\_plc.py <{}<{}'EOF'}\dots\texttt{EOF}):
\begin{lstlisting}[language=Python]
from pymodbus.server import StartTcpServer
from pymodbus.datastore import (
    ModbusSequentialDataBlock, ModbusSlaveContext, ModbusServerContext,
)

# Fake "process": registers 0..9 lie about a sensor that reads 4242.
block = ModbusSequentialDataBlock(0, [4242] * 100)
slave = ModbusSlaveContext(hr=block, ir=block, co=block, di=block)
ctx   = ModbusServerContext(slaves=slave, single=True)

print("Fake PLC listening on 0.0.0.0:502 -- registers report 4242")
StartTcpServer(context=ctx, address=("0.0.0.0", 502))
\end{lstlisting}
Run it:
\begin{lstlisting}[language=bash]
python3 /tmp/fake_plc.py
\end{lstlisting}
\end{stepbox}

\begin{stepbox}{6.3 -- Confirm the spoof from a second shell}
\begin{lstlisting}[language=bash]
docker compose exec student \
  python3 /labs/scripts/modbus_read.py
\end{lstlisting}
\emph{Expected output:}
\begin{lstlisting}
Registers 0..9: [4242, 4242, 4242, 4242, 4242,
                 4242, 4242, 4242, 4242, 4242]
\end{lstlisting}
The client cannot tell. If an HMI was pointed at this IP, it would now
display \texttt{4242} as if it were a real pump pressure. Stop the fake
server with Ctrl-C and bring the real PLC back:
\begin{lstlisting}[language=bash]
docker compose start openplc
\end{lstlisting}
\end{stepbox}

\section*{7 -- Cleanup}

\begin{stepbox}{7.1 -- Tear down}
\begin{lstlisting}[language=bash]
exit                          # leave the student shell
docker compose down           # keep volumes; -v also wipes them
ls bachelor/labs/_stack/captures/
\end{lstlisting}
Your three PCAPs (\texttt{lab04-mb.pcap}, \texttt{lab04-ua.pcap},
\texttt{lab04-mq.pcap}) should sit on the host -- they are mounted from
the container.
\end{stepbox}

\begin{deliverbox}
Hand in:
\begin{itemize}[leftmargin=*, nosep]
  \item All three PCAPs (Modbus, OPC UA, MQTT).
  \item Annotated MBAP byte table (Step 2.3) -- mark every byte with its
        meaning and the IT analogue you would compare it to.
  \item One screenshot of the OpenPLC \emph{Monitoring} page showing
        \texttt{QX0.0 = TRUE} after Step 2.4.
  \item The tshark excerpt that shows \texttt{SecurityPolicy=None} and
        \texttt{UserIdentityToken=Anonymous} from Step 3.3.
  \item A wildcard-subscribe transcript from Step 4.4.
  \item The filled-in side-by-side table from Step 5.1.
  \item A one-paragraph answer: ``\emph{Which of the three protocols
        could you put on a hostile network today, with what extra
        controls, and why?}''
  \item (Stretch) Source of your fake PLC plus a screenshot of the
        client receiving \texttt{4242}.
\end{itemize}
\end{deliverbox}

\begin{warnbox}
Lab containers only. The injection, the wildcard subscribe, and the
fake PLC all work because the lab broker, PLC, and OPC UA server were
deliberately deployed with default insecure settings. Running any of
these against equipment you do not own is, depending on jurisdiction,
either a felony or a billable safety incident -- and often both.
\end{warnbox}
\end{document}
